\section{Supporting Nondeterministism}\label{sec:extension}

The core language of Section~\ref{sec:context-typing} is
deterministic: every term reduces to at most one value under any
environment.  This section extends the language with
\emph{nondeterminism}, expressed in terms of primitive generators that
can produce multiple values across different executions, and
identifies the adjustments required to the type denotation.

Concretely, our language now supports generators over primitive types:
\begin{align*}
  \primop ::= \quad & \ldots ~|~ \Code{bool\_gen()} ~|~
  \Code{nat\_gen()} ~|~ \ldots
\end{align*}
The nondeterministic choice $e_1 \oplus e_2$ is then expressible as
$\zif{\Code{bool\_gen()}}{e_1}{e_2}$.  These generators serve as the
``impure external environments'' introduced informally in
Section~\ref{sec:context-typing}: the environment $\Envir_x(e) \doteq
\zlet{x}{\Code{nat\_gen()}}{e}$ is now a first-class term that
introduces a nondeterministic binding for $x$ whose capability is
\emph{independent} of the surrounding context at each invocation.

Nondeterminism raises two technical challenges.  First, nondeterministic
generators can be invoked multiple times, producing independent
capabilities at each call; the type system must account for this
\emph{resource duplication}.  Second, nondeterminism breaks the
single-valued assumption underlying our type denotation, requiring a
generalization of the denotation of function and sum types.
% (\S\ref{subsec:nondet-denotation}).  We conclude with case studies
% (\S\ref{subsec:case-studies}) demonstrating the expressiveness of the
% resulting system.

\subsection{Resource Duplication and the Demonic Modality}\label{subsec:duplication}

In the deterministic core language, every function application produces
a result whose capability is dependent on the function's capability via
the default dependent conjunction ($\ctcomma$) in \textsc{T-Let}.  Consider two
successive calls to a deterministic function $f$:
\begin{align*}
  &\underline{f{:}\urt{\Unit}{\top} \ctsarr \urt{\Nat}{\top}
    \ctcomma x{:}\urt{\Nat}{\top} \ctcomma y{:}\urt{\Nat}{\top}
    \vdash e : \tau}
  \\&f{:}\urt{\Unit}{\top} \ctsarr \urt{\Nat}{\top} \vdash
    \zlet{x}{f\;()}{\zlet{y}{f\;()}{e}} : \tau
\end{align*}
Both $x$ and $y$ are entangled with $f$'s capability: knowing $f$
constrains both, so their capabilities cannot be treated as
independent.  This is correct for a deterministic $f$, whose output
is fully determined by its input.

With a nondeterministic generator like $\Code{nat\_gen()}$, however,
successive calls produce genuinely independent values.  The key
observation is that non-deterministic generators like
$\Code{nat\_gen()}$ must be given a \emph{demonic} parameter type:
\begin{align*}
  \Code{nat\_gen} : \ort{\Unit}{\top} \ctsarr \urt{\Nat}{\top}
\end{align*}
The demonic parameter type $\ort{\Unit}{\top}$ asserts that the
generator's result is independent of its argument in a specific sense:
since $\ort{\Unit}{\top}$ is a safety (overapproximate) type, it
requires the argument to be consistent with $\top$, i.e., any unit
value, but places no angelic obligation on it.  Crucially, a demonic
type in a \emph{negative position} (as a function parameter) behaves
like a persistent resource in the context algebra: it can be freely
duplicated across multiple uses without entangling the corresponding
capabilities.

This follows from a key algebraic property of the context logic: the
demonic modality $\demonic P$ is persistent by definition, i.e., once
$\demonic P$ holds, it continues to hold as more environments are
added to the capability. This means that a demonic capability, once
established, remains available regardless of what other environments
are added to the context, which is exactly the duplication behavior
needed for resource generators.

%% To see
%% why, recall that $r \models \demonic P$ $\exists r' \supseteq r$
%% s.t. $r' \models P$.  Now suppose r⊨\demonicPr \models \demonic P
%% r⊨\demonicP and consider any extension r′′r'' r′′ of rr r in the
%% set-inclusion sense, i.e., r⊆r′′r \subseteq r'' r⊆r′′. Since r′⊇rr'
%% \supseteq r r′⊇r and r′′⊇rr'' \supseteq r r′′⊇r, we can take r′∪r′′r'
%% \cup r'' r′∪r′′ as a witness: it satisfies r′∪r′′⊇r′′r' \cup r''
%% \supseteq r'' r′∪r′′⊇r′′ and r′∪r′′⊇r′⊨Pr' \cup r'' \supseteq r'
%% \models P r′∪r′′⊇r′⊨P by monotonicity of PP P under superset
%% extension. Thus r′′⊨\demonicPr'' \models \demonic P r′′⊨\demonicP,
%% showing that the demonic modality is preserved under set-inclusion
%% extension — i.e., 

Concretely, this means that two successive calls to $\Code{nat\_gen()}$
can be typed with independent capabilities for $x$ and $y$:
\begin{align*}
  &\underline{\Code{nat\_gen}{:}\ort{\Unit}{\top} \ctsarr \urt{\Nat}{\top}
    \ctstar x{:}\urt{\Nat}{\top} \ctstar y{:}\urt{\Nat}{\top}
    \vdash e : \tau}
  \\&\underline{\Code{nat\_gen}{:}\ort{\Unit}{\top} \ctsarr \urt{\Nat}{\top}
    \ctcomma x{:}\urt{\Nat}{\top} \ctcomma y{:}\urt{\Nat}{\top}
    \vdash e : \tau}
  \\&\Code{nat\_gen}{:}\ort{\Unit}{\top} \ctsarr \urt{\Nat}{\top}
    \vdash \zlet{x}{\Code{nat\_gen}\;()}{\zlet{y}{\Code{nat\_gen}\;()}{e}} : \tau
\end{align*}
The key step is the shift from dependent ($\ctcomma$) to disjoint
($\ctstar$) conjunction in the second-to-first line.  This shift is
justified by the demonic modality's persistence: since
$\Code{nat\_gen}$'s parameter type is demonic, its capability can be
duplicated freely, allowing the two results $x$ and $y$ to be treated
as disjoint.  No separate persistency modality is needed — the
persistence arises directly from the demonic character of the parameter
type.

More generally, any function whose parameter type is demonic
(overapproximate) can be invoked multiple times with independent result
capabilities, since the demonic parameter type's persistence in negative
position enables the necessary resource duplication.  Functions with
angelic parameter types, by contrast, require their argument to cover
a specific set of values, which entangles successive calls via the
shared argument capability.

\begin{example}[Coverage Types]\label{ex:coverage}
  The coverage type system of~\citet{CoverageType} verifies that a
  test generator produces every desirable input.  Consider the
  following BST generator:
\begin{lstlisting}
let rec bst_gen (r: nat) (lo: nat)
                (range: nat -> nat -> nat): nat tree =
  if r <= 0 then Leaf
  else Leaf $\oplus$
    let r' = range (1, r) in
    let left  = bst_gen (r' - 1) lo range in
    let right = bst_gen (r - r') (lo + r') range in
    Node r' left right
\end{lstlisting}
  The coverage specification asserts that for any valid range oracle,
  \texttt{bst\_gen} produces every BST whose keys fall in $(\Code{lo},
  \Code{lo}+r]$:
  \begin{align*}
    &\Code{val bst\_gen} :
      r{:}\ort{\Nat}{\top} \ctsarr lo{:}\ort{\Nat}{\top} \ctsarr
    \\&\quad\quad
      \Code{range}{:}\underbrace{(x{:}\ort{\Nat}{\top} \ctsarr
        y{:}\ort{\Nat}{x < \vnu} \ctsarr
        \urt{\Nat}{x < \vnu < y})}_{\text{demonic parameters; angelic result}}
      \ctwand
    \\&\quad\quad
      \urt{\Nat\;\Code{tree}}{\Code{bst}(\vnu) \land
        \forall x.\, \Code{mem}(\vnu, x) \iff
        (\Code{lo}+1 \leq x \leq \Code{lo}+r)}
  \end{align*}
  The range oracle has demonic parameter types ($\ort{\Nat}{\top}$ and
  $\ort{\Nat}{x < \vnu}$) and an angelic result type
  ($\urt{\Nat}{x < \vnu < y}$).  The demonic parameters mean that
  successive calls to \texttt{range} with different arguments can be
  typed with independent result capabilities — the oracle's output at
  each recursive call is disjoint from its outputs at other calls.
  This is precisely the resource duplication property established
  above: the demonic parameter types make the oracle's capability
  persistent in negative position, enabling the recursive calls to
  accumulate independent angelic capabilities for the left and right
  subtrees.  The $\ctwand$ on the oracle ensures it is disjoint from
  the size $r$ and lower-bound $\Code{lo}$, so the oracle's choices
  do not depend on those parameters.  The angelic result type then
  asserts that every valid BST is reachable by some sequence of
  oracle choices.
\end{example}

\subsection{Type Denotation for Nondeterminism}\label{subsec:nondet-denotation}

Nondeterminism breaks the single-valued assumption of the deterministic
language in two ways that require adjustments to the type denotation.

\paragraph{Function Types}
In the deterministic setting, a function term $e$ reduces to a unique
value, and the denotation of a function type can reason about that
value directly.  With nondeterminism, $e$ can reduce to multiple
different function values, and different values may behave differently
when applied.  Consider:
\begin{align*}
  f_1 \doteq \zlam{x}{\Unit}{\true \oplus \false}
  \qquad
  f_2 \doteq (\zlam{x}{\Unit}{\true}) \oplus (\zlam{x}{\Unit}{\false})
\end{align*}
Both $f_1\;()$ and $f_2\;()$ can reduce to either $\true$ or $\false$.
However, they behave differently in a higher-order context.  Consider:
\begin{align*}
  \emptyset \vdash \zzlam{f}{f\;() == f\;()} :
  f{:}(\ort{\Unit}{\top} \ctsarr \urt{\Bool}{\top}) \ctwand
  \urt{\Bool}{\top}
\end{align*}
This function applies $f$ twice and checks equality.  With $f_1$,
each call independently produces $\true$ or $\false$, so the result
covers both.  With $f_2$, call-by-value semantics first evaluates
$f_2$ to either $\zlam{x}{\Unit}{\true}$ or $\zlam{x}{\Unit}{\false}$
before binding it to $f$; both subsequent calls then return the same
value, so the result is always $\true$.  The original denotation
cannot distinguish these two cases, since it does not account for
which function value $e$ reduces to.

\paragraph{Sum Types}
A similar issue arises for sum types.  The judgement:
\begin{align*}
  \emptyset \vdash \Code{bool\_gen()} :
  \ort{\Bool}{\vnu = \true} \ctoplus \ort{\Bool}{\vnu = \false}
\end{align*}
should hold, since $\Code{bool\_gen()}$ nondeterministically produces
either $\true$ or $\false$.  However, the original denotation
$\denot{\tau_1 \ctoplus \tau_2}\;e \doteq \denot{\tau_1}\;e \choplus
\denot{\tau_2}\;e$ requires the \emph{context} to be split across the
two branches, not the term itself.  Under the empty context, there is
no capability to split, so the judgement fails despite being
semantically valid.

\paragraph{Generalized Denotation}
Both issues stem from the same root cause: the original denotation
evaluates $e$ once and reasons about the result, but a
nondeterministic $e$ can reduce to different values in different
executions.  The fix is to universally quantify over all possible
reduction results of $e$, lifting $e$'s nondeterminism into the
type context:
\begin{align*}
  \denot{x{:}\tau_x \ctsarr \tau}_\Sigma\;e
    &\doteq \forall f.\; \denot{f \doteq e}_\Sigma \chimpl
             \denot{\tau_x}_\Sigma\;x \chimpl
             \denot{\tau}_{\Sigma,x{:}\eraserf{\tau_x}}\;(f\;x)
  \\\denot{x{:}\tau_x \ctwand \tau}_\Sigma\;e
    &\doteq \forall f.\; \denot{f \doteq e}_\Sigma \chimpl
             \denot{\tau_x}_\Sigma\;x \chwand
             \denot{\tau}_{\Sigma,x{:}\eraserf{\tau_x}}\;(f\;x)
  \\\denot{\tau_1 \ctoplus \tau_2}_\Sigma\;e
    &\doteq \forall x.\; \denot{x \doteq e}_\Sigma \chimpl
             \denot{\tau_1}_\Sigma\;x \choplus \denot{\tau_2}_\Sigma\;x
\end{align*}
In each case, a universally quantified fresh variable ($f$ or $x$)
ranges over all possible reduction results of $e$, and the original
denotation is applied to that variable.  In the deterministic case,
$e$ has a unique reduction result and the quantifier is trivially
eliminated, so the generalized denotation coincides with the original.

For the function-type issue, the generalized denotation correctly
distinguishes $f_1$ and $f_2$.  The type
$\ort{\Unit}{\top} \ctsarr \urt{\Bool}{\top}$ now requires that for
every function value $f$ that $e$ can reduce to, every boolean is
reachable by applying $f$ to $()$.  Since $f_2$ reduces to a
\emph{specific} constant function, not every boolean is reachable from
a single reduction result, so $f_2$ correctly fails to have this type:
\begin{align*}
  &\denot{\ort{\Unit}{\top} \ctsarr \urt{\Bool}{\top}}\;
    (\zlam{x}{\Unit}{\true}) \oplus (\zlam{x}{\Unit}{\false})
  \\=\;& \forall f.\; \Code{Atom}(f = \zlam{x}{\Unit}{\true} \lor
          f = \zlam{x}{\Unit}{\false}) \chimpl
          \angelic\,\Code{Atom}(f\;() = \true \lor f\;() = \false)
  \quad\judgfail
\end{align*}
For the sum-type issue, the generalized denotation lifts
$\Code{bool\_gen()}$'s nondeterminism into the context via the fresh
variable $x$, allowing the sum operator to partition it across the
two branches as usual:
\begin{align*}
  &\denot{\ort{\Bool}{\vnu = \true} \ctoplus
          \ort{\Bool}{\vnu = \false}}\;\Code{bool\_gen()}
  \\=\;& \forall x.\; \Code{Atom}(x = \true \lor x = \false) \chimpl
          \demonic\,\Code{Atom}(\vnu = \true) \choplus
          \demonic\,\Code{Atom}(\vnu = \false)
  \quad\judgok
\end{align*}

\subsection{Case Studies}\label{subsec:case-studies}

We illustrate the expressiveness of context types with two further
case studies beyond the BST generator of Example~\ref{ex:coverage}:
a reverse reachability property for list reversal, and a disprove
reasoning task for a merge function.

\subsubsection{Reverse Reachability}\label{subsubsec:reverse}

The standard list reversal function admits two distinct reachability
specifications:
\begin{lstlisting}
let rec reverse (l: list nat) : list nat =
  match l with
  | []      -> []
  | x :: xs -> reverse xs ++ [x]
\end{lstlisting}
\begin{align*}
  &\Code{val reverse} :
    l{:}\ort{\Nat\;\Code{list}}{\top} \ctsarr
    \urt{\Nat\;\Code{list}}{\Code{rev}(\vnu, l)}
  \\&\Code{val reverse} :
    l{:}\urt{\Nat\;\Code{list}}{\Code{sortedIncr}(\vnu)} \ctsarr
    \urt{\Nat\;\Code{list}}{\Code{sortedDecr}(\vnu)}
\end{align*}
The first uses a demonic input type: for any list $l$,
\texttt{reverse} produces its reversal.  The second uses an angelic
input type: when $l$ ranges over all increasingly-sorted lists, the
output ranges over all decreasingly-sorted lists.  The second
specification is strictly stronger, demonstrating that context types
can express reachability properties that depend on the full range of
the input.

\subsubsection{Disprove Reasoning}\label{subsubsec:disprove}

The following function merges two lists element-by-element:
\begin{lstlisting}
let rec merge (l1: elem list) (l2: elem list) : elem list =
  match l1, l2 with
  | l1, []          -> l1
  | [], l2          -> l2
  | h1::t1, h2::t2 ->
      let t = merge t1 t2 in
      if h1 < h2 then h1 :: h2 :: t
      else if h1 > h2 then h2 :: h1 :: t
      else h1 :: t
\end{lstlisting}
Context types can express both positive reachability specifications
and \emph{disprove} reasoning — showing that a function cannot always
produce a sorted output.  The following two specifications hold:
\begin{align*}
  &\Code{val merge} :
    l_1{:}\ort{\Nat\;\Code{list}}{\top} \ctsarr
    l_2{:}\urt{\Nat\;\Code{list}}{\top} \ctwand
    (\ort{\Nat\;\Code{list}}{\neg\Code{sortedIncr}(\vnu)}
      \ctoplus \ort{\Nat\;\Code{list}}{\Code{sortedIncr}(\vnu)})
  \\&\Code{val merge} :
    l_1{:}\urt{\Nat\;\Code{list}}{\top} \ctwand
    l_2{:}\urt{\Nat\;\Code{list}}{\top} \ctwand
    (\ort{\Nat\;\Code{list}}{\neg\Code{sortedIncr}(\vnu)}
      \ctoplus \ort{\Nat\;\Code{list}}{\top})
\end{align*}
The first says: when $l_2$ ranges over all lists, the output either
fails to be sorted or is sorted.  The second strengthens this: when
both inputs range over all lists, the output can fail to be sorted.
The following specification correctly fails to type-check, confirming
that \texttt{merge} cannot be specified as always producing a sorted
output:
\begin{align*}
  &\Code{val merge} :
    l_1{:}\ort{\Nat\;\Code{list}}{\top} \ctsarr
    l_2{:}\ort{\Nat\;\Code{list}}{\top} \ctsarr
    \ort{\Nat\;\Code{list}}{\Code{sortedIncr}(\vnu)}
  \quad\judgfail
\end{align*}
This failed judgement serves as a reachability-based witness to the
function's incorrectness: there exist inputs for which
\texttt{merge} does not produce a sorted list.
